\documentclass[12pt,a4paper,fontset=none]{ctexart}

\usepackage{fontspec}
\usepackage{xeCJK}
\setmainfont{Times New Roman}
\setCJKmainfont{SimSun}[BoldFont=SimHei]

\usepackage{geometry}
\geometry{a4paper, margin=2.5cm}

\usepackage{amsmath,amssymb,amsthm}
\usepackage{xcolor}
\usepackage{mdframed}
\usepackage{etoolbox}
\usepackage{fancyhdr}
\usepackage{titlesec}
\usepackage{hyperref}

\usepackage{tocloft}
\renewcommand{\cftsecleader}{\hfill}
\renewcommand{\cftsubsecleader}{\cftdotfill{\cftdotsep}}

\definecolor{capellaBlue}{RGB}{0,102,204}
\definecolor{capellaRed}{RGB}{204,0,0}
\definecolor{capellaGreen}{RGB}{0,204,0}
\definecolor{capellaYellow}{RGB}{204,153,0}

\mdfdefinestyle{definitionstyle}{
    linecolor=capellaBlue,
    leftline=true,
    rightline=false,
    topline=false,
    bottomline=false,
    backgroundcolor=capellaBlue!6,
    linewidth=2pt,
    innertopmargin=6pt,
    innerbottommargin=6pt,
    skipabove=4pt,
    skipbelow=4pt
}

\mdfdefinestyle{exercisestyle}{
    linecolor=capellaYellow,
    leftline=true,
    rightline=false,
    topline=false,
    bottomline=false,
    backgroundcolor=capellaYellow!6,
    linewidth=2pt,
    innertopmargin=6pt,
    innerbottommargin=6pt,
    skipabove=4pt,
    skipbelow=4pt
}

\mdfdefinestyle{examplestyle}{
    linecolor=capellaRed,
    leftline=true,
    rightline=false,
    topline=false,
    bottomline=false,
    backgroundcolor=capellaRed!6,
    linewidth=2pt,
    innertopmargin=6pt,
    innerbottommargin=6pt,
    skipabove=4pt,
    skipbelow=4pt
}

\mdfdefinestyle{theoremstyle}{
    linecolor=capellaGreen,
    leftline=true,
    rightline=false,
    topline=false,
    bottomline=false,
    backgroundcolor=capellaGreen!6,
    linewidth=2pt,
    innertopmargin=6pt,
    innerbottommargin=6pt,
    skipabove=4pt,
    skipbelow=4pt
}

\newtheoremstyle{bluehead}
  {0pt}{0pt}
  {\normalfont}
  {}
  {\bfseries\color{capellaBlue}}
  {}
  { }
  {\thmname{#1}\ \thmnumber{#2}\ \thmnote{\ {\textcolor{capellaBlue}{(#3)}}}}

% red heading for examples
\newtheoremstyle{redhead}
  {0pt}{0pt}
  {\normalfont}
  {}
  {\bfseries\color{capellaRed}}
  {}
  { }
  {\thmname{#1}\ \thmnumber{#2}\ \thmnote{\ {\textcolor{capellaRed}{(#3)}}}}

\newtheoremstyle{greenhead}
  {0pt}{0pt}
  {\normalfont}
  {}
  {\bfseries\color{capellaGreen}}
  {}
  { }
  {\thmname{#1}\ \thmnumber{#2}\ \thmnote{\ {\textcolor{capellaGreen}{(#3)}}}}

% yellow heading for exercises
\newtheoremstyle{yellowhead}
  {0pt}{0pt}
  {\normalfont}
  {}
  {\bfseries\color{capellaYellow}}
  {}
  { }
  {\thmname{#1}\ \thmnumber{#2}\ \thmnote{\ {\textcolor{capellaYellow}{(#3)}}}}

\newcounter{theorem_counter}

% definitions in blue
\theoremstyle{bluehead}
\newtheorem{definition}{定义}[subsection]
\surroundwithmdframed[style=definitionstyle]{definition}

% theorems in blue (same formatting as definitions)
\theoremstyle{greenhead}
\newtheorem{theorem}{定理}[subsection]
\surroundwithmdframed[style=theoremstyle]{theorem}

% exercises in yellow
\theoremstyle{yellowhead}
\newtheorem{exercise}{习题}[subsection]
\surroundwithmdframed[style=exercisestyle]{exercise}

% examples in red (same formatting except color)
\theoremstyle{redhead}
\newtheorem{example}{例题}[subsection]
\surroundwithmdframed[style=examplestyle]{example}


\pagestyle{fancy}
\fancyhf{}
\fancyhead[L]{ CHM NOTE}
\fancyhead[R]{\leftmark}
\fancyfoot[C]{\thepage}


\begin{document}

\begin{titlepage}
\centering
\vspace*{\fill}
\begin{minipage}{\textwidth}
  \centering
  {\Huge\bfseries 预科二下半学期数学笔记\par}
  \vspace{1cm}
  {\Large 堡堡菲\par}
  \vspace{1cm}
    {\large \today\par}
\end{minipage}
\vspace*{\fill}
\end{titlepage}


\tableofcontents
\thispagestyle{empty}
\newpage

\section*{导言}
\addcontentsline{toc}{section}{导言}

我说把 设$\varepsilon>0$ 写成 $\forall \varepsilon>0$ 的入平时成绩零分, 你们尔朵龙吗?

\newpage


\section{数列极限}


\subsection{Stolz 定理}

\begin{theorem}
  (Stolz 定理).

  \hspace{-2em}设 $\{x_n\}$ 和 $\{y_n\}$ 是两串实数序列, 如果 $\{y_n\}$ 满足以下条件中的任何一个:

  \vspace{0.1cm}
  
  1) $\{y_n\}$ 严格单调增且 $y_n \to +\infty$;

  \vspace{0.1cm}
  
  2) $\{y_n\}$ 严格单调减且 $x_n,\ y_n \to 0$.

  \vspace{0.2cm}
  
  \hspace{-2em}那么 $\displaystyle \lim_{n \to \infty} \frac{x_{n+1}-x_n}{y_{n+1}-y_n}=L \Rightarrow \lim_{n \to \infty} \frac{x_n}{y_n}=L$ (其中 $L$ 为有限或$\pm \infty$).
\end{theorem}

\hspace{-2em}\textbf{证} \quad 先证明情况 1). 只对有限的 $L$ 作证明. 改写条件为
$$
\text{设} \varepsilon>0, \exists N \in \mathbb{N}, \forall n>N, \left|\frac{x_{n}-x_{n+1}}{y_{n}-y_{n+1}}-L\right|<\varepsilon.
$$
结合 $\{y_n\}$ 的单调性条件自然地去掉绝对值, 得到
$$
(L-\varepsilon)(y_{n}-y_{n+1})<x_{n}-x_{n+1}<(L+\varepsilon)(y_{n}-y_{n+1}).
$$
任取 $m>n$, 将上述不等式中的 $n$ 换成 $n+1, n+2, \cdots, m-1$ 然后将所有不等式相加 
$$
(L-\varepsilon)(y_{n}-y_{m})<x_{n}-x_{m}<(L+\varepsilon)(y_{n}-y_{m}).
$$
令 $m \to \infty$, $x_m, y_m \to 0$, 得到.
$$
\left|\frac{x_n}{y_n}-L\right|<\varepsilon.
$$

再证明情况 2). 只对有限的 $L$ 作证明. 设 $\varepsilon>0$, $\exists N \in \mathbb{N}$, $\forall n \ge N$, 则
$$(L-\varepsilon)(y_{n+1}-y_n)<b_{n+1}-b_n<(L+\varepsilon)(y_{n+1}-y_n).$$
仿照上面的过程,再次进行累加, 并重新整理为绝对值
$$
\left|\frac{x_{n}-x_{N}}{y_{n}-y_{N}}-L\right|<\varepsilon.
$$
为了得到关于 $\displaystyle \left|\frac{x_n}{y_n}-L\right|$ 的估计, 我们注意到恒等式
$$
\frac{x_n}{y_n}-L =  \left(1-\frac{y_N}{y_n}\right) \left(\frac{x_n-x_N}{y_n-y_N}-L\right) + \frac{x_N-Ly_N}{y_n}.
$$
再取一个无穷指标 $N_1$, 使得当 $n > N_1$ 时成立
$$
0<1-\frac{y_N}{y_n}<2, \quad \left|\frac{x_N-Ly_N}{y_n}\right|<\varepsilon.
$$
则在 $n > \max \{N,N_1\}$ 时就得到
\begin{equation*}
\left|\frac{x_n}{y_n}-L\right|<3\varepsilon.\tag*{\qed}
\end{equation*}

\begin{example}
  计算数列极限 
  $$
  \lim_{n \to \infty} \left(\frac{2}{2^2-1}\right)^{\frac{1}{2^{n-1}}} \left(\frac{2^2}{2^3-1}\right)^{\frac{1}{2^{n-2}}} \cdots \left(\frac{2^{n-1}}{2^n-1}\right)^{\frac{1}{2}}.
  $$
\end{example}

\hspace{-2em}\textbf{解} \quad 利用对数函数的连续性, 将极限转化为计算级数 
$$
\sum_{k=1}^{n-1} \frac{1}{2^{n-k}} \ln \left( \frac{2^k}{2^{k+1}-1} \right) = \frac{1}{2^{n-1}}\sum_{k=1}^{n-1} 2^{k-1} \ln \left(\frac{2^k}{2^{k+1}-1}\right).
$$
注意到
$$
\lim_{n \to \infty} \frac{2^{n-2}\ln\left(\frac{2^{n-1}}{2^n-1}\right)}{2^{n-1}-2^{n-2}}=\ln \left( \frac{1}{2} \right).
$$
由 Stolz 定理, 上述级数的极限为 $\displaystyle \ln \left( \frac{1}{2} \right)$, 从而原数列的极限为 $\displaystyle \frac{1}{2}$.\hfill\(\qed\)

\vspace{0.2cm}

\begin{exercise}
已知数列 $\{x_n\}_{n \ge 1}$ 的初值为 $1$, 并满足如下的递推关系
$$
x_{n+1}=x_n+3\sqrt{x_n}+\frac{n}{\sqrt{x_n}}.
$$
证明数列 $\displaystyle \left\{\frac{x_n}{n^2}\right\}$ 有极限, 并求出这个极限.
\end{exercise}

\begin{theorem}
  (Cesàro 定理).
  设 $\{a_n\}_{n \ge 1}$ 为实数序列, 记 $\displaystyle \sigma_n=\frac{a_1+a_2+\cdots +a_n}{n}$. 那么 $$\displaystyle \lim_{n \to \infty} a_n=a \Rightarrow \displaystyle \lim_{n \to \infty} \sigma_n=a.$$
\end{theorem}

Cesàro 定理是 Stolz 定理的一个特例.

\vspace{0.3cm}

\begin{example}
  已知 $\displaystyle \lim_{n \to \infty} x_n=a$, $\displaystyle \lim_{n \to \infty} y_n=b$, 求证
  $$
  \lim_{n \to \infty} \frac{x_1y_n+x_2y_{n-1}+\cdots+x_{n-1}y_2+x_ny_1}{n} = ab.
  $$
\end{example}

\hspace{-2em}\textbf{证} \quad 注意到(这里 $a$ 的变形需要引入一个指标 $N$ 以保证)
$$
\left|\frac{\sum_{i=1}^nx_iy_{n+1-i}}{n}-ab\right|=\left|\frac{\sum_{i=1}^nx_iy_{n+1-i}}{n}-\left(\frac{\sum_{i=1}^nx_i}{n}\right)b\right|=\left|\frac{\sum_{i=1}^nx_i(y_{n+1-i}-b)}{n}\right|.
$$
后续证法与标准的 Cesàro 定理的证明完全相同, 这里不再赘述.\hfill\(\qed\)

\vspace{0.2cm}

\begin{exercise}

(参数条件见定理1.2.2)

1) Cesàro 定理的逆命题是否成立?

2) 对 $k \ge 1$, 记 $b_k=a_{k+1}-a_{k}$. 证明, $\forall n \in \mathbb{N}_{\ge 2}$, 都有 $\displaystyle a_n-\sigma_n=\frac{1}{n}\sum_{k=1}^{n-1}kb_k$.

3) 设 $\{kb_k\}_{k \ge 1}$ 有界并且 $\{\sigma_{n}\}_{n \ge 1}$ 收敛, 证明 $\{a_n\}$ 也收敛.
\end{exercise}

\begin{theorem}
  (一般情况下的 Stolz 定理).

  \hspace{-2em}设 $\{x_n\}$, $\{y_n\}$ 是两个实数列, 且 $\{y_n\}$ 单调无界, 那么
  $$
  \liminf_{n \to \infty} \frac{x_{n+1}-x_n}{y_{n+1}-y_n} \le \liminf_{n \to \infty} \frac{x_n}{y_n} \le \limsup_{n \to \infty} \frac{x_n}{y_n} \le \limsup_{n \to \infty} \frac{x_{n+1}-x_n}{y_{n+1}-y_n}.
  $$
\end{theorem}

\subsection{其他例题}

\begin{example}
  设 $|a|>1$, $\alpha \in \mathbb{R^+}$, 那么 $\displaystyle \lim_{n \to \infty}\frac{n^\alpha}{a^n}=0$.
\end{example}

\hspace{-2em} \textbf{证} \quad 将 $a^n$ 形式上记为 $(1+(a-1))^n$ 并利用二项式定理展开, 保留第 $[\alpha]+1$ 项即可. \hfill\(\qed\)

\vspace{0.2cm}

\begin{exercise}
  证明恒等式: $\displaystyle C_{n}^{k-1}+C_{n}^{k}=C_{n+1}^k$.
\end{exercise}

\newpage

\section{第二章}

\begin{definition}
目录会自动显示本节标题。
\end{definition}

\end{document}











